% =====================================================================
%  Préambule commun aux documents de soutenance
% =====================================================================
\usepackage[T1]{fontenc}
\usepackage[utf8]{inputenc}
\usepackage[sfdefault]{carlito}
\usepackage[french]{babel}
\usepackage{csquotes}
\usepackage[a4paper,margin=2.2cm,top=2.5cm,bottom=2.2cm]{geometry}
\usepackage{microtype}
\usepackage[htt]{hyphenat}
\tolerance=2000
\emergencystretch=3em
\widowpenalty=10000
\clubpenalty=10000
\setlength{\parskip}{5pt plus 2pt minus 1pt}

\usepackage{xcolor}
\usepackage{colortbl}
\usepackage{array}
\usepackage{tabularx}
\usepackage{longtable}
\usepackage{enumitem}
\usepackage{ragged2e}
\usepackage{titlesec}
\usepackage{fancyhdr}
\usepackage{tcolorbox}
\usepackage[hidelinks]{hyperref}

\definecolor{btkblue}{HTML}{14507A}
\definecolor{btklight}{HTML}{1B6CA8}
\definecolor{btkpale}{HTML}{EAF3FA}
\definecolor{btkambre}{HTML}{9A6B12}
\definecolor{btkambrepale}{HTML}{FDF6E6}
\definecolor{btkvert}{HTML}{2E7D5B}
\definecolor{btkvertpale}{HTML}{EDF7F2}

\renewcommand{\tabularxcolumn}[1]{>{\RaggedRight\arraybackslash}p{#1}}
\newcolumntype{L}[1]{>{\RaggedRight\arraybackslash}p{#1}}
\renewcommand{\arraystretch}{1.25}

\titleformat{\section}{\Large\bfseries\color{btkblue}}{\thesection}{0.7em}{}
\titleformat{\subsection}{\large\bfseries\color{btklight}}{\thesubsection}{0.6em}{}
\titlespacing*{\section}{0pt}{2.4ex plus 1ex}{1.2ex}
\titlespacing*{\subsection}{0pt}{1.8ex plus .8ex}{0.8ex}

\pagestyle{fancy}
\fancyhf{}
\renewcommand{\headrulewidth}{0.4pt}
\renewcommand{\headrule}{\hbox to\headwidth{\color{btklight}\leaders\hrule height \headrulewidth\hfill}}
\fancyhead[R]{\small\color{btkblue}Oussema Korbosli --- PFE BTK}
\fancyfoot[C]{\small\thepage}

% Encadré « à savoir par cœur »
\newtcolorbox{coeur}[1][]{
  colback=btkpale, colframe=btklight, boxrule=0.9pt, arc=2pt,
  left=7pt, right=7pt, top=5pt, bottom=5pt,
  fonttitle=\bfseries\small, title={À savoir par c\oe{}ur}, #1}

% Encadré « note » : conseil de présentation
\newtcolorbox{note}[1][]{
  colback=btkvertpale, colframe=btkvert, boxrule=0.8pt, arc=2pt,
  left=7pt, right=7pt, top=4pt, bottom=4pt, #1}

% Encadré « chiffre corrigé »
\newtcolorbox{corrige}[1][]{
  colback=btkambrepale, colframe=btkambre, boxrule=0.9pt, arc=2pt,
  left=7pt, right=7pt, top=5pt, bottom=5pt,
  fonttitle=\bfseries\small, title={Chiffre corrigé}, #1}

% Question / réponse
\newcommand{\question}[1]{\subsection{#1}}
\newcommand{\cle}[1]{\par\vspace{2pt}\noindent
  \colorbox{btkpale}{\parbox{\dimexpr\linewidth-2\fboxsep}{%
  \textbf{Phrase à dire :} \emph{#1}}}\par\vspace{3pt}}
